%%%
%%% Radical "⽼"
%%%
\section*{Radical 125: ``⽼'' (耂)}\addcontentsline{toc}{section}{Radical 125: ⽼,耂}\addcontentsline{loh}{figure}{\#\#\#\# 125: ⽼}

%%%%%%%%%% 老 %%%%%%%%%%
\subsection*{老}\addcontentsline{loh}{figure}{老}

\begin{Entry}{老}{6}{⽼}[Kangxi 125]
  \begin{Phonetics}{老}{lao3}[][HSK 1,2]
    \definition*{s.}{Sobrenome: Lao}
    \definition{adj.}{velho; envelhecido; idade avançada | antigo; de longa data; existe há muito tempo | antigo; desatualizado; obsoleto; ultrapassado  | antigo; tradicional; original | coberto de vegetação; os vegetais cresceram além do período ideal para serem consumidos | resistente; endurecido; alimentos muito cozidos | (sobre cores) escuro; profundo | último nascido; o mais novo | veterano; experiente; sofisticado}
    \definition{adv.}{longo; por muito tempo | sempre (fazendo algo) | muito}
    \definition{pref.}{usado para designar pessoas, ordem de classificação, certos nomes de animais e plantas}
    \definition{s.}{idosos; pessoas mais velhas | ancião; sênior; um título respeitoso para pessoas mais velhas}
    \definition{v.}{envelhecer | morrer; referindo"-se à morte de um idoso}
  \seealsoref{旧}{jiu4}
  \synonymref{陈}{chen2}
  \synonymref{古}{gu3}
  \synonymref{故}{gu4}
  \synonymref{练}{lian4}
  \synonymref{逝}{shi4}
  \synonymref{熟}{shu2}
  \synonymref{远}{yuan3}
  \synonymref{长}{zhang3}
  \antonymref{存}{cun2}
  \antonymref{近}{jin4}
  \antonymref{嫩}{nen4}
  \antonymref{少}{shao4}
  \antonymref{生}{sheng1}
  \antonymref{首}{shou3}
  \antonymref{鲜}{xian1}
  \antonymref{小}{xiao3}
  \antonymref{新}{xin1}
  \antonymref{幼}{you4}
  \end{Phonetics}
\end{Entry}

\begin{Entry}{老人}{6,2}{⽼,⼈}
  \begin{Phonetics}{老人}{lao3ren2}[][HSK 1]
    \definition[位,名,个]{s.}{homem ou mulher de idade avançada; o idoso; o velho}
  \end{Phonetics}
\end{Entry}

\begin{Entry}{老人家}{6,2,10}{⽼,⼈,⼧}
  \begin{Phonetics}{老人家}{lao3ren5jia5}[][HSK 7-9]
    \definition[位,名,个]{s.}{avô; avó; pessoa idosa venerável; um título respeitoso para os idosos | maneira de chamar o pai ou a mãe idosos na frente dos outros; referir"-se aos próprios pais ou aos pais, professores, etc. de outras pessoas}
  \end{Phonetics}
\end{Entry}

\begin{Entry}{老乡}{6,3}{⽼,⼄}
  \begin{Phonetics}{老乡}{lao3xiang1}[][HSK 6]
    \definition[个,位]{s.}{conterrâneo; conterrâneo | uma maneira de chamar um fazendeiro cujo nome você não conhece}
  \synonymref{农民}{nong2min2}
  \end{Phonetics}
\end{Entry}

\begin{Entry}{老大}{6,3}{⽼,⼤}
  \begin{Phonetics}{老大}{lao3da4}[][HSK 7-9]
    \definition{adj.}{velho; mais velho}
    \definition{adv.}{muito; extremamente (encontrado com mais frequência no chinês vernáculo antigo)}
    \definition{s.}{filho mais velho; o mais velho entre os irmãos | capitão de um barco; mestre de uma embarcação à vela; o barqueiro principal do barco de madeira, e também um termo genérico para barqueiros | líder; chefe; o título que algumas gangues ou grupos do crime organizado usam para seus líderes}
  \end{Phonetics}
\end{Entry}

\begin{Entry}{老公}{6,4}{⽼,⼋}
  \begin{Phonetics}{老公}{lao3gong1}
    \definition[个,位,名]{s.}{marido; esposo}
  \antonymref{男人}{nan2ren2}
  \antonymref{男人}{nan2ren5}
  \antonymref{丈夫}{zhang4fu5}
  \end{Phonetics}
  \begin{Phonetics}{老公}{lao3gong5}[][HSK 4]
    \definition[个,位,名]{s.}{marido; esposo}
  \synonymref{男人}{nan2ren2}
  \synonymref{男人}{nan2ren2}
  \synonymref{丈夫}{zhang4fu5}
  \end{Phonetics}
\end{Entry}

\begin{Entry}{老化}{6,4}{⽼,⼔}
  \begin{Phonetics}{老化}{lao3hua4}[][HSK 7-9]
    \definition{s.}{Química: envelhecer | (pessoas) envelhecer; ficar velho | (conhecimento, etc.) tornam"-se obsoletos}
    \definition{s.}{envelhecimento; período de amadurecimento; pré"-queima; maturação}
  \end{Phonetics}
\end{Entry}

\begin{Entry}{老太太}{6,4,4}{⽼,⼤,⼤}
  \begin{Phonetics}{老太太}{lao3tai4tai5}[][HSK 3]
    \definition[位,名,个]{s.}{velha senhora; (em tratamento direto) Venerável Senhora; uma maneira respeitosa de chamar uma senhora idosa; título honorífico para mulheres idosas | (forma de tratamento) sua velha mãe; minha velha mãe, avó ou sogra; referindo"-se à própria mãe, à mãe do outro ou à mãe de outra pessoa, à sogra ou à sogra política}
  \synonymref{母亲}{mu3qin5}
  \synonymref{岳母}{yue4mu3}
  \end{Phonetics}
\end{Entry}

\begin{Entry}{老头儿}{6,5,2}{⽼,⼤,⼉}
  \begin{Phonetics}{老头儿}{lao3tou2r5}[][HSK 3]
    \definition{s.}{(coloquial) (com um tom de intimidade) velho; velho amigo}
  \seealsoref{老头子}{lao3tou2zi5}
  \end{Phonetics}
\end{Entry}

\begin{Entry}{老头子}{6,5,3}{⽼,⼤,⼦}
  \begin{Phonetics}{老头子}{lao3tou2zi5}
    \definition{s.}{velho antiquado (ou velho rabugento) | (referindo"-se ao marido idoso) meu velho | chefe de uma sociedade secreta | Coloquial: velho; velho rabugento}
  \seealsoref{老头儿}{lao3tou2r5}
  \end{Phonetics}
\end{Entry}

\begin{Entry}{老汉}{6,5}{⽼,⽔}
  \begin{Phonetics}{老汉}{lao3han4}[][HSK 7-9]
    \definition{pron.}{(usado por um homem idoso) eu, eu mesmo}
    \definition[个,位]{s.}{homem velho; homens idosos; um velho | um velho como eu; utilizado para autorreferência}
  \end{Phonetics}
\end{Entry}

\begin{Entry}{老字号}{6,6,5}{⽼,⼦,⼝}
  \begin{Phonetics}{老字号}{lao3zi4hao5}[][HSK 7-9]
    \definition[家]{s.}{nome antigo no ramo comercial; empresa (ou loja) de longa data; empreendimento (ou loja) antigo e famoso | loja, empresa ou marca de mercadorias com reputação consolidada; lojas tradicionais}
  \end{Phonetics}
\end{Entry}

\begin{Entry}{老师}{6,6}{⽼,⼱}
  \begin{Phonetics}{老师}{lao3shi1}[][HSK 1]
    \definition[个,位,名,些]{s.}{professor; título honorífico para professores; refere"-se, de maneira geral, a pessoas que transmitem cultura e tecnologia ou que são dignas de admiração em termos de ideias, moralidade e conhecimentos profissionais}
  \seealsoref{先生}{xian1sheng5}
  \synonymref{教练}{jiao4lian4}
  \synonymref{教师}{jiao4shi1}
  \synonymref{教授}{jiao4shou4}
  \synonymref{师长}{shi1zhang3}
  \synonymref{师傅}{shi1fu5}
  \antonymref{弟子}{di4zi3}
  \antonymref{徒弟}{tu2di5}
  \antonymref{学生}{xue2sheng5}
  \end{Phonetics}
\end{Entry}

\begin{Entry}{老年}{6,6}{⽼,⼲}
  \begin{Phonetics}{老年}{lao3nian2}[][HSK 2]
    \definition[个]{s.}{idoso; velhice; idade acima de 60 ou 70 anos}
  \synonymref{晚年}{wan3nian2}
  \antonymref{青年}{qing1nian2}
  \antonymref{少年}{shao4nian2}
  \antonymref{童年}{tong2nian2}
  \antonymref{中年}{zhong1nian2}
  \end{Phonetics}
\end{Entry}

\begin{Entry}{老百姓}{6,6,8}{⽼,⽩,⼥}
  \begin{Phonetics}{老百姓}{lao3bai3xing4}[][HSK 3]
    \definition[些]{s.}{povo; civis; pessoas comuns; residentes (em contraste com militares e funcionários públicos)}
  \seealsoref{人民}{ren2min2}
  \synonymref{公民}{gong1min2}
  \synonymref{民众}{min2zhong4}
  \synonymref{平民}{ping2min2}
  \synonymref{群众}{qun2zhong4}
  \antonymref{干部}{gan4bu4}
  \antonymref{官员}{guan1yuan2}
  \antonymref{军人}{jun1ren2}
  \antonymref{领导}{ling3dao3}
  \end{Phonetics}
\end{Entry}

\begin{Entry}{老伴}{6,7}{⽼,⼈}
  \begin{Phonetics}{老伴}{lao3ban4}
    \definition[个]{s.}{marido ou esposa (de um casal de idosos)}
  \seealsoref{老伴儿}{lao3ban4r5}
  \end{Phonetics}
\end{Entry}

\begin{Entry}{老伴儿}{6,7,2}{⽼,⼈,⼉}
  \begin{Phonetics}{老伴儿}{lao3ban4r5}[][HSK 7-9]
    \definition{s.}{marido ou esposa}
  \seealsoref{老伴}{lao3ban4}
  \end{Phonetics}
\end{Entry}

\begin{Entry}{老兵}{6,7}{⽼,⼋}
  \begin{Phonetics}{老兵}{lao3bing1}
    \definition{s.}{velho soldado; veterano de guerra | veterano (alguém que tem muita experiência em alguma área)}
  \end{Phonetics}
\end{Entry}

\begin{Entry}{老远}{6,7}{⽼,⾡}
  \begin{Phonetics}{老远}{lao3yuan3}[][HSK 7-9]
    \definition{adj.}{muito longe; longínquo; de grande distância}
  \end{Phonetics}
\end{Entry}

\begin{Entry}{老实}{6,8}{⽼,⼧}
  \begin{Phonetics}{老实}{lao3shi5}[][HSK 4]
    \definition{adj.}{franco; sincero; honesto | bom; bem"-comportado | ingênuo; simplório; meio bobo; facilmente enganado; eufemismo para pouco inteligente}
  \seealsoref{诚实}{cheng2shi5}
  \seealsoref{实在}{shi2zai4}
  \synonymref{诚恳}{cheng2ken3}
  \synonymref{诚挚}{cheng2zhi4}
  \synonymref{敦厚}{dun1hou4}
  \synonymref{厚道}{hou4dao5}
  \synonymref{忠实}{zhong1shi2}
  \antonymref{奸诈}{jian1zha4}
  \antonymref{狡猾}{jiao3hua2}
  \antonymref{淘气}{tao2//qi4}
  \antonymref{调皮}{tiao2pi2}
  \antonymref{顽皮}{wan2pi2}
  \antonymref{虚伪}{xu1wei3}
  \end{Phonetics}
\end{Entry}

\begin{Entry}{老实说}{6,8,9}{⽼,⼧,⾔}
  \begin{Phonetics}{老实说}{lao3shi5shuo1}[][HSK 7-9]
    \definition{expr.}{``Honestamente.''; ser honesto; ser franco; falando francamente; sinceramente, muitas vezes funciona como uma interjeição em uma frase para indicar ênfase}
  \end{Phonetics}
\end{Entry}

\begin{Entry}{老朋友}{6,8,4}{⽼,⽉,⼜}
  \begin{Phonetics}{老朋友}{lao3peng2you3}[][HSK 2]
    \definition[个,位,名]{s.}{velho amigo; refere"-se a amigos que conhecemos há muito tempo e com quem temos uma relação íntima}
  \end{Phonetics}
\end{Entry}

\begin{Entry}{老板}{6,8}{⽼,⽊}
  \begin{Phonetics}{老板}{lao3ban3}[][HSK 3]
    \definition[个,位]{s.}{chefe; dono; líder; refere"-se ao gerente de uma empresa comercial ou industrial | antigo título honorífico dado a atores famosos de ópera ou atores que também eram diretores de companhias de ópera}
  \synonymref{店主}{dian4zhu3}
  \synonymref{雇主}{gu4zhu3}
  \antonymref{伙计}{huo3ji5}
  \end{Phonetics}
\end{Entry}

\begin{Entry}{老虎}{6,8}{⽼,⾌}
  \begin{Phonetics}{老虎}{lao3hu3}
    \definition[只]{s.}{tigre | instalações que consomem grande quantidade de energia ou matérias"-primas, ou que geram resíduos ou danos significativos; metaforicamente, refere"-se a máquinas e equipamentos que consomem grandes quantidades de energia ou matérias"-primas | pessoa muito gananciosa; pessoa que comete corrupção, roubo ou sonegação fiscal em grande escala; essa metáfora descreve alguém que usa seu poder para cometer crimes em larga escala, como peculato, roubo e sonegação fiscal}
  \seealsoref{虎}{hu3}
  \end{Phonetics}
\end{Entry}

\begin{Entry}{老是}{6,9}{⽼,⽇}
  \begin{Phonetics}{老是}{lao3shi4}[][HSK 2]
    \definition{adv.}{sempre; indica que a ação continua ou que o estado permanece inalterado, equivalente a 一直}
  \seealsoref{一直}{yi4zhi2}
  \synonymref{常常}{chang2chang2}
  \synonymref{经常}{jing1chang2}
  \synonymref{频频}{pin2pin2}
  \synonymref{总是}{zong3shi4}
  \antonymref{偶尔}{ou3'er3}
  \antonymref{有时}{you3shi2}
  \end{Phonetics}
\end{Entry}

\begin{Entry}{老家}{6,10}{⽼,⼧}
  \begin{Phonetics}{老家}{lao3jia1}[][HSK 4]
    \definition{s.}{cidade natal; local de origem | lugar nativo; refere"-se às gerações anteriores da família ou ao local onde a pessoa nasceu ou viveu}
  \synonymref{故乡}{gu4xiang1}
  \synonymref{家乡}{jia1xiang1}
  \synonymref{家园}{jia1yuan2}
  \synonymref{祖籍}{zu3ji2}
  \antonymref{外地}{wai4di4}
  \end{Phonetics}
\end{Entry}

\begin{Entry}{老婆}{6,11}{⽼,⼥}
  \begin{Phonetics}{老婆}{lao3po5}[][HSK 4]
    \definition[个,位,名]{s.}{esposa}
  \seealsoref{妻子}{qi1zi5}
  \synonymref{夫人}{fu1ren5}
  \synonymref{太太}{tai4tai5}
  \synonymref{媳妇}{xi2fu5}
  \antonymref{老公}{lao3gong1}
  \antonymref{老汉}{lao3han4}
  \antonymref{老公}{lao3gong5}
  \end{Phonetics}
\end{Entry}

%%%%%%%%%% 考 %%%%%%%%%%
\subsection*{考}\addcontentsline{loh}{figure}{考}

\begin{Entry}{考}{6}{⽼}
  \begin{Phonetics}{考}{kao3}[][HSK 1]
    \definition*{s.}{Sobrenome: Kao}
    \definition{adj.}{antigo; velho; com idade avançada}
    \definition{s.}{o pai falecido de alguém}
    \definition{v.}{examinar; dar (fazer) um exame, teste ou questionário | verificar; inspecionar | estudar; verificar; investigar | perguntar; testar; fazer perguntas para que o outro responda, a fim de testar suas habilidades em determinada área}
  \seealsoref{考试}{kao3//shi4}
  \synonymref{查}{cha2}
  \synonymref{检}{jian3}
  \synonymref{究}{jiu1}
  \synonymref{问}{wen4}
  \end{Phonetics}
\end{Entry}

\begin{Entry}{考生}{6,5}{⽼,⽣}
  \begin{Phonetics}{考生}{kao3sheng1}[][HSK 2]
    \definition{s.}{candidato a exame; alunos inscritos para o exame de admissão}
  \end{Phonetics}
\end{Entry}

\begin{Entry}{考场}{6,6}{⽼,⼟}
  \begin{Phonetics}{考场}{kao3chang3}[][HSK 6]
    \definition{s.}{sala de exames}
  \synonymref{反对}{fan3dui4}
  \synonymref{拒绝}{ju4jue2}
  \synonymref{抗拒}{kang4ju4}
  \end{Phonetics}
\end{Entry}

\begin{Entry}{考试}{6,8}{⽼,⾔}
  \begin{Phonetics}{考试}{kao3//shi4}[][HSK 1]
    \definition[场,个,次]{s.}{teste; exame; prova; atividades realizadas para verificar conhecimentos ou habilidades}
    \definition{v.+compl.}{testar; avaliar; avaliar conhecimentos e habilidades por meio de perguntas escritas ou orais.}
  \seealsoref{考}{kao3}
  \synonymref{测试}{ce4shi4}
  \synonymref{测验}{ce4yan4}
  \synonymref{考察}{kao3cha2}
  \synonymref{考核}{kao3he2}
  \synonymref{试验}{shi4yan4}
  \end{Phonetics}
\end{Entry}

\begin{Entry}{考核}{6,10}{⽼,⽊}
  \begin{Phonetics}{考核}{kao3he2}[][HSK 5]
    \definition{v.}{examinar; checar; avaliar; avaliar (a proficiência de alguém)}
  \synonymref{考试}{kao3//shi4}
  \synonymref{考察}{kao3cha2}
  \synonymref{审核}{shen3he2}
  \end{Phonetics}
\end{Entry}

\begin{Entry}{考虑}{6,10}{⽼,⾌}
  \begin{Phonetics}{考虑}{kao3lv4}[][HSK 4]
    \definition{v.}{considerar; refletir sobre; levar em conta}
  \seealsoref{打算}{da3suan5}
  \synonymref{商量}{shang1liang5}
  \synonymref{思考}{si1kao3}
  \synonymref{思索}{si1suo3}
  \synonymref{探讨}{tan4tao3}
  \synonymref{推敲}{tui1qiao1}
  \synonymref{研究}{yan2jiu1}
  \synonymref{研讨}{yan2tao3}
  \synonymref{琢磨}{zuo2mo5}
  \antonymref{鲁莽}{lu3mang3}
  \antonymref{武断}{wu3duan4}
  \end{Phonetics}
\end{Entry}

\begin{Entry}{考验}{6,10}{⽼,⾺}
  \begin{Phonetics}{考验}{kao3yan4}[][HSK 3]
    \definition[场,个,种]{s.}{teste; julgamento; atividade realizada para verificar se as habilidades, ideias, moral e qualidades de uma pessoa atendem aos requisitos}
    \definition{v.}{testar; testar as capacidades, ideias, moral e qualidades de uma pessoa através de situações, ações ou ambientes difíceis, para verificar se elas atendem aos requisitos}
  \synonymref{检验}{jian3yan4}
  \synonymref{验证}{yan4zheng4}
  \end{Phonetics}
\end{Entry}

\begin{Entry}{考量}{6,12}{⽼,⾥}
  \begin{Phonetics}{考量}{kao3liang2}[][HSK 7-9]
    \definition{v.}{considerar; examinar e medir}
  \end{Phonetics}
\end{Entry}

\begin{Entry}{考察}{6,14}{⽼,⼧}
  \begin{Phonetics}{考察}{kao3cha2}[][HSK 4]
    \definition{v.}{inspecionar; investigar; observar e estudar}
  \synonymref{参观}{can1guan1}
  \synonymref{调查}{diao4cha2}
  \synonymref{观察}{guan1cha2}
  \synonymref{考核}{kao3he2}
  \synonymref{视察}{shi4cha2}
  \end{Phonetics}
\end{Entry}

\begin{Entry}{考题}{6,15}{⽼,⾴}
  \begin{Phonetics}{考题}{kao3ti2}[][HSK 6]
    \definition{s.}{questões de exame; prova de exame; tópicos de exame}
  \synonymref{试题}{shi4ti2}
  \synonymref{题目}{ti2mu4}
  \end{Phonetics}
\end{Entry}

%%%%%%%%%% 者 %%%%%%%%%%
\subsection*{者}\addcontentsline{loh}{figure}{者}

\begin{Entry}{者}{8}{⽼}
  \begin{Phonetics}{者}{zhe3}[][HSK 3]
    \definition*{s.}{Sobrenome: Zhe}
    \definition{part.}{significa 是 e é usado após palavras, frases e orações para indicar uma pausa}
    \definition{pron.}{usado para se referir à pessoa, coisa ou assunto que realiza uma ação ou possui um determinado atributo | pessoas; caras (Usado para se referir a alguém envolvido em uma determinada profissão, que acredita em uma determinada ideologia ou que tem uma forte tendência para algo) | usado após certos números ou palavras direcionais para se referir a coisas mencionadas anteriormente | significado semelhante a 这 (mais comum na linguagem coloquial antiga)}
  \seealsoref{是}{shi4}
  \seealsoref{这}{zhe4}
  \end{Phonetics}
\end{Entry}

%%%%% EOF %%%%%

